\chapter{El ruido cuántico}
\label{cap:ruido}

% GUION. Fuente principal: doc/03_Ruido_Cuantico.md (desarrollado ago-2026).
% Cifras medidas sobre nuestro chip: doc/memoria/materiales.md seccion 2.

\section{Estados mixtos y matriz de densidad}
\label{sec:densidad}
Por qué un vector de estado deja de bastar cuando hay ruido: el sistema pasa a ser una
\textbf{mezcla estadística}. Definición de $\rho$ y cálculo del valor esperado como
$\mathrm{Tr}(\rho O)$.

\textbf{Consecuencia práctica que se usa en el capítulo~\ref{cap:rigor}:} el coste de memoria
pasa de $2^n$ a $4^n$. Para 15 qubits son 17,2\,GB por proceso.

\section{Canales cuánticos y operadores de Kraus}
\label{sec:kraus}
Formalismo $\mathcal{E}(\rho) = \sum_k E_k \rho E_k^\dagger$ y por qué permite componer el
ruido puerta a puerta.

\textbf{Sutileza relevante:} el ruido \textbf{no conmuta} con la unitaria. Dos descomposiciones
equivalentes del mismo circuito acumulan decoherencia de forma distinta —lo que explica un
resultado del capítulo~\ref{cap:rigor}.

\section{Las tres fuentes de error}
\label{sec:fuentes}

\subsection{Decoherencia: T1 y T2}
Relajación de energía y pérdida de coherencia de fase. La cota física $T_2 \leq 2T_1$.
Lo que importa no es $T_1$ en absoluto, sino \textbf{cuánto dura el circuito comparado con
$T_1$}.

% Fuente: src/features.py, familia F
\textbf{Y la duración no es la suma de las puertas:} las que actúan sobre qubits distintos se
ejecutan simultáneamente. Lo que cuenta es el \textbf{camino crítico del DAG}, entre 2 y 3
veces menor que la suma.

\subsection{Errores de puerta}
Imperfección de los pulsos. Dos hallazgos que solo se ven mirando datos reales:
\begin{itemize}
    \item la \texttt{rz} tiene error \textbf{exactamente cero} porque es \textbf{virtual}
          —cambia la fase de referencia y no emite pulso—;
    \item \texttt{id}, \texttt{sx} y \texttt{x} traen valores \textbf{idénticos}, porque IBM
          las deriva del mismo pulso calibrado.
\end{itemize}

\subsection{Errores de lectura}
Confusión entre $\ket{0}$ y $\ket{1}$ al medir. \textbf{Es asimétrico}, porque relajarse es
más fácil que excitarse.

\section{El modelo de ruido de este trabajo}
\label{sec:modelo_ruido}
% Fuente: src/quantum_gen.py, HardwareTelemetrics
Los tres canales canónicos de Qiskit Aer, \textbf{calibrados contra el error medido de cada
puerta física de cada día}. Validación: la infidelidad efectiva del modelo iguala el
\texttt{gate\_error} que reporta IBM.

\subsection{Qué no se modela: el crosstalk}
% Fuente: doc/03_Ruido_Cuantico.md seccion 4
Los tres canales son \textbf{puramente locales}. Declararlo como limitación, y contar el
experimento del grado de conectividad: se propuso como \textit{feature}, se midió
($\rho = +0{,}025$, $p = 0{,}62$) y se descartó, porque \textbf{sin crosstalk no existe camino
causal} entre el grado del qubit y la etiqueta.

\section{Mitigar frente a corregir}
\label{sec:qem_vs_qec}
QEC necesita qubits que no hay; QEM estima el error y lo descuenta. Dónde se sitúa este
trabajo: mitigación, y en la variante \textbf{pre-ejecución}.
